\chapter{Memory: Reconstitution and the Recursive Limit}
\label{ch:memory}

\begin{quotation}
\noindent\textit{Synopsis.} This chapter returns to memory as a concrete
illustration of reconstitution. Memories are analyzed as traces requiring
reconstruction through schemas, priors, and contextual cues; eyewitness
testimony and autobiographical recall are examined as continuation under
severe informational constraint. The chapter shows that memory does not
escape the framework's central limitations but reproduces them in
concentrated form, making visible, at the scale of a single mind, dynamics
that are otherwise easiest to see in large institutional relay systems.
\end{quotation}

This chapter proceeds under the narrow reading of MEM|8 identified in
Chapter~\ref{ch:mem8}; wherever a claim below would be affected by the
broad reading instead, this is noted explicitly rather than left
implicit.

\section{Memory as a Three-Stage Relay}

Standard treatments of memory in cognitive psychology distinguish
encoding, consolidation, and retrieval as separate processes, each with
its own resource constraints. This structure maps directly onto
\cref{def:substrate-chain}.

\begin{construction}[Memory as a Substrate Chain]
\label{con:memory-relay}
Let $A_0$ be an experienced event, $A_1$ the encoded trace immediately
following it, $A_2$ the consolidated trace some time later, and $A_3$
the recollection produced at retrieval. The transition $A_0 \to A_1$ is
governed by encoding resources (attention, salience) as its budget
$\beta_0$; $A_1 \to A_2$ by consolidation resources (rehearsal, sleep,
interference from subsequent events) as $\beta_1$; and $A_2 \to A_3$ by
retrieval resources (available cues, time pressure) as $\beta_2$.
\end{construction}

Each transition in \cref{con:memory-relay} satisfies
\cref{def:budget-relay}'s independence condition in the ordinary case: an
event's encoding is not typically decided with reference to what a much
later retrieval attempt will need, consolidation proceeds according to
its own resource constraints without reference to encoding's particular
priorities, and a specific retrieval attempt draws on whatever cues are
available at that moment without renegotiating what consolidation
preserved. \cref{con:memory-relay} therefore is not merely analogous to a
budget relay, in the way Chapter~\ref{ch:fidelity-continuum} was careful
to establish for reconstitution generally (\cref{prop:recursive-insufficiency}):
it is a three-stage instance of one, with each stage's budget
independently and plausibly set.

\begin{proposition}[Memory Retrieval Is a Relay Transition]
\label{prop:memory-relay}
The transition $\text{trace} \to \text{recollection}$ satisfies the
conditions of a substrate-chain transition (\cref{def:substrate-chain}).
\end{proposition}

\begin{proofsketch}
Retrieval takes a residual trace and a current schema and produces a
usable recollection, incurring costs (effort, plausibility constraints
supplied by the schema) and losses (aspects of the original event the
trace does not constrain) in the sense of \cref{def:budget-discard}. This
satisfies \cref{def:substrate-chain} directly, and is the final stage of
\cref{con:memory-relay}.
\end{proofsketch}

\section{Schema-Driven Confident Error}

Eyewitness-memory research documents confident, richly reconstructed
testimony that is nonetheless factually wrong: the schema fills a gap the
trace does not constrain, and the result is presented, to the rememberer
and to others, as though it were retrieved rather than regenerated. This
is a concrete instance of \cref{prop:recursive-insufficiency}'s ``priors
patch'' (\cref{sec:coord-priors}) operating at retrieval time on an
already-insufficient trace, and is structurally the same failure mode as
\cref{prop:open-causal-space}'s painting case, transposed from an
external image to an internal reconstruction. In terms of
\cref{con:memory-relay}, this is exactly the pattern
\cref{thm:composition-failure} predicts: a distinction lost or degraded
at encoding or consolidation -- because $\tau_0$ or $\tau_1$, whatever
was salient or well-rehearsed at the time, did not prioritize it -- is
unavailable to retrieval regardless of how confidently retrieval's own
schema fills the resulting gap, since confidence at $A_3$ cannot restore
what $A_1$ or $A_2$ never carried forward.

\begin{proposition}[Memory Does Not Escape Composition Failure]
\label{prop:memory-no-escape}
Memory's reconstitutive character does not constitute an exception to
\cref{thm:composition-failure}.
\end{proposition}

\begin{proofsketch}
By \cref{prop:memory-relay} and \cref{con:memory-relay}, memory retrieval
is a relay transition, indeed the final stage of a three-stage relay; by
\cref{prop:recursive-insufficiency}, this transition inherits the
open-task-space vulnerability of any other stage in a budget relay.
\cref{thm:composition-failure} therefore applies to memory exactly as it
applies to the media-pipeline case of Chapter~\ref{ch:pipelines},
contrary to an informal suggestion, made earlier in the research program
that produced this monograph, that reconstitution somehow ``changes the
game'' rather than relocating the same problem to a new substrate.
\end{proofsketch}
